\input{../tex-inputs/latex-leseansicht-vorspann}

\section[Arthur und Olga Schnitzler an Gustav Schwarzkopf, 16. 5. 1904]{L04415 Arthur und Olga Schnitzler an Gustav Schwarzkopf, 16. 5. 1904}
\nopagebreak\mylabel{L04415v}
\rehead{ }\normalsize\beginnumbering\briefempfaengerindex{Schwarzkopf, Gustav@\textsc{Schwarzkopf, Gustav}!zzzSchnitzler, Olga@\emph{von Olga Schnitzler}!1904-05-163@{16. 5. 1904}|(be}\briefempfaengerindex{Schwarzkopf, Gustav@\textsc{Schwarzkopf, Gustav}!zzzSchnitzler, Arthur@\emph{von Arthur Schnitzler}!1904-05-163@{16. 5. 1904}|(be}
\toendnotes[C]{\smallbreak\pagebreak[2]}
\correspDesc{Versand  durch Arthur Schnitzler, Olga Schnitzler am 16. 5. 1904 in Monreale
\newline{}Übermittlung  am 17. 5. 1904 in Palermo
\newline{}Zustellung  am 19. 5. 1904 in Wien 1/1 1
\newline{}Erhalt  durch Gustav Schwarzkopf am 19. 5. 1904 in Tiefer Graben 23}\toendnotes[C]{\smallbreak}
\Standort{DLA, A:Schnitzler, HS.1985.1.1897.}
\physDesc{Bildpostkarte, 269 Zeichen
\newline{}Handschrift Arthur Schnitzler: Bleistift, deutsche Kurrent
\newline{}Handschrift Olga Schnitzler: Bleistift, lateinische Kurrent
\newline{}Versand: 1) Stempel: »\nobreak{}\oindex{Palermo@\textbf{Palermo}|pwk}Palermo Ferr\textcolor{gray}{o}{[}via{]}, 1\textcolor{gray}{7 5 04}\nobreak{}«.   2) Stempel: »\nobreak{}\oindex{I., Innere Stadt@\textbf{I., Innere Stadt}|pwk}Wien 1/\textcolor{gray}{1} 1, 19./\textcolor{gray}{5. 04}, 8–9½V., Bestellt\nobreak{}«. }\toendnotes[C]{\smallbreak}\pstart{}{\pb}\textsc{M Gustav Schwarzkopf}\pend{}\pstart{}\textsc{Vienne I}\oindex{Wien@\textbf{Wien}|pw}\pend{}\pstart{}\textsc{Tiefer Graben 23}\oindex{Wien@\textbf{Wien}!I., Innere Stadt@\textbf{I., Innere Stadt}!Tiefer Graben 23@\textbf{Tiefer Graben 23}|pw}.\pend{}\pstart{}\textsc{Austria}\oindex{Österreich@\textbf{Österreich}|pw}.\pend{}{\bigskip}
\pstart
           \centering{}{\pb}\textcolor{gray}{\textbf{MONREALE}}\oindex{Monreale@\textbf{Monreale}|pw}\pend
           
\pstart
           \centering{}\textcolor{gray}{\textbf{Convento dei Benedettini\oindex{Chiostro di Monreale@\textbf{Chiostro di Monreale}|pw} Veduta del Chiostro\oindex{Chiostro di Monreale@\textbf{Chiostro di Monreale}|pw} col Campanile – secolo XII.}}\pend
           \vspace{1em}
\pstart
           {\pb}16. 5.\pend
           \vspace{0.5em}
\pstart
           Viele Grüße u{ }ſchönen Dank für die lieben \label{K_L04415-1v}\edtext{\substVorne{}\textsuperscript{Ka}\substDazwischen{}Na\substHinten{}chrichten}{\lemma{\textnormal{\emph{Nachrichten}}}\Cendnote{\textnormal{XXXX Auszeichnungsfehler: Dokument L04313 nicht gefunden.}}}\label{K_L04415-1}. Wann{ }ſoll die \textsc{Première}\pwindex{Schwarzkopf, Max 12.\,6.\,1857 Wien – 14.\,4.\,1928 ebd.@\textsc{Schwarzkopf, Max} (12.\,6.\,1857 Wien – 14.\,4.\,1928 ebd.)!reine Tor. Gesellschaftsstück in vier Akten@\strich\emph{Der reine Tor. Gesellschaftsstück in vier Akten}|pwv}\eventindex{Neues Deutsches Theater@\textbf{Neues Deutsches Theater}!Uraufführung von Der reine Tor, 17.6.1904@Uraufführung von Der reine Tor, 17.6.1904|pwv}{ }ſein?\pend
           
\pstart
           Nachricht nach Taormina\oindex{Taormina@\textbf{Taormina}|pw}{ }\label{K_L04415-2v}\edtext{\textsc{postrest}}{\lemma{\textnormal{\emph{postrest}}}\Cendnote{\textnormal{poste restante: postlagernd}}}\label{K_L04415-2} erbeten. –\pend
           
\pstart
           Herzlich{\\[\baselineskip]}Ihr \spacefill\mbox{A.}\pend
           \leftskip=0em{}\selectlanguage{ngerman}\vspace{1em}
\pstart
           \noindent{}{[}hs. O. Schnitzler:{]} Herzl. Grüsse, und schreiben Sie bald wieder so einen lieben Brief!\pend
           \pstart \spacefill\mbox{Olga S.}\pend{}\selectlanguage{ngerman}\endnumbering\briefempfaengerindex{Schwarzkopf, Gustav@\textsc{Schwarzkopf, Gustav}!zzzSchnitzler, Olga@\emph{von Olga Schnitzler}!1904-05-163@{16. 5. 1904}|)be}\briefempfaengerindex{Schwarzkopf, Gustav@\textsc{Schwarzkopf, Gustav}!zzzSchnitzler, Arthur@\emph{von Arthur Schnitzler}!1904-05-163@{16. 5. 1904}|)be}\mylabel{L04415h}
\begin{anhang}
\end{anhang}\newcommand{\dateiname}{L04415}\newcommand{\titel}{Arthur und Olga Schnitzler an Gustav Schwarzkopf, 16. 5. 1904}\newcommand{\editorInnen}{Herausgegeben von Jahnke, SelmaMüller, Martin Anton}\input{../tex-inputs/latex-leseansicht-abspann}
